\subsection{Threat Modeling dengan STRIDE}

Setelah aset, \textit{entry points}, dan \textit{trust boundaries} pada sistem DecMed berhasil diidentifikasi, langkah selanjutnya adalah melakukan analisis ancaman menggunakan metode STRIDE. Metode ini digunakan untuk mengidentifikasi berbagai ancaman keamanan yang berpotensi memengaruhi komponen sistem berdasarkan enam kategori ancaman, yaitu \textit{Spoofing}, \textit{Tampering}, \textit{Repudiation}, \textit{Information Disclosure}, \textit{Denial of Service}, dan \textit{Elevation of Privilege}.

Analisis dilakukan terhadap seluruh elemen pada DFD, yang mencakup entitas eksternal, proses, penyimpanan data, serta aliran data. Untuk setiap komponen, diidentifikasi kemungkinan skenario serangan yang dapat terjadi beserta dampak yang ditimbulkannya terhadap operasional dan keamanan sistem. Hasil identifikasi ancaman ini selanjutnya digunakan untuk menentukan informasi dan bukti digital yang perlu dicatat oleh sistem. Dengan demikian, kebutuhan \textit{audit trail} yang dirumuskan pada penelitian ini dapat ditelusuri secara langsung berdasarkan ancaman yang berpotensi terjadi pada DecMed.

Hasil identifikasi ancaman menggunakan metode STRIDE pada komponen-komponen sistem DecMed ditunjukkan pada Tabel~\ref{tab:stride-all-components}.

\begin{xltabular}{\textwidth}{|c|>{\raggedright\arraybackslash}p{3.5cm}|c|X|X|}
    \caption{Analisis STRIDE Komponen Sistem DecMed}
    \label{tab:stride-all-components} \\
    \hline
    \textbf{ID} & \textbf{Komponen (Kode)} & \textbf{STRIDE} & \textbf{Skenario} & \textbf{Dampak} \\
    \hline
    \endfirsthead

    \multicolumn{5}{c}{{Tabel \thetable\ -- Lanjutan dari halaman sebelumnya}} \\
    \hline
    \textbf{ID} & \textbf{Komponen (Kode)} & \textbf{STRIDE} & \textbf{Skenario Singkat} & \textbf{Dampak} \\
    \hline
    \endhead

    \hline
    \multicolumn{5}{r}{{\textit{Berlanjut ke halaman berikutnya...}}} \\
    \endfoot

    \hline
    \endlastfoot

    T01 & Entitas Manusia (E1, E2, E3) & S & Penyamaran menggunakan PIN/kredensial curian. & Akses EMR ilegal. \\
    \hline
    T02 & Entitas Manusia (E1, E2, E3) & R & Menyangkal pernah mengizinkan atau mengubah data. & Sengketa medis/hukum. \\
    \hline
    T03 & Entitas Sistem Eksternal (E4, E5) & S & Node luar menyamar untuk memalsukan status jaringan. & Gagal eksekusi transaksi. \\
    \hline
    T04 & Entitas Sistem Eksternal (E4, E5) & R & Node menyangkal menerima pesanan transaksi gas/data. & Status \textit{on-chain} terputus. \\
    \hline
    T05 & Proses Klien (P1, P2, P3) & S & \textit{Malware} menyamar sebagai proses aplikasi sah. & \textit{Bypass} kontrol lokal. \\
    \hline
    T06 & Proses Klien (P1, P2, P3) & T & Modifikasi \textit{runtime} memori aplikasi di perangkat. & Eksekusi kode jahat. \\
    \hline
    T07 & Proses Klien (P1, P2, P3) & R & Aplikasi gagal/dicegah mencatat log lokal. & Forensik lokal hilang. \\
    \hline
    T08 & Proses Klien (P1, P2, P3) & I & \textit{Memory dump} membocorkan Kunci IOTA/PIN. & Kredensial terkompromi. \\
    \hline
    T09 & Proses Klien (P1, P2, P3) & D & Menguras CPU/RAM perangkat klien. & Aplikasi \textit{crash}/hang. \\
    \hline
    T10 & Proses Klien (P1, P2, P3) & E & Mengakali UI untuk mengakses menu admin. & Akses fitur terlarang. \\
    \hline
    T11 & Proses \textit{Backend} (P4, P5, P6, P7, P8) & S & Injeksi JWT bajakan ke \textit{Proxy API}. & Sesi pengguna diambil alih. \\
    \hline
    T12 & Proses \textit{Backend} (P4, P5, P6, P7, P8) & T & Manipulasi instruksi memori di \textit{PRE Engine}. & Enkripsi data EMR korup. \\
    \hline
    T13 & Proses \textit{Backend} (P4, P5, P6, P7, P8) & R & Gagal mengirim metadata persetujuan ke IOTA. & \textit{Audit trail} terputus. \\
    \hline
    T14 & Proses \textit{Backend} (P4, P5, P6, P7, P8) & I & Pesan \textit{error/stack-trace} membocorkan kFrag/Token. & Data rahasia terekspos. \\
    \hline
    T15 & Proses \textit{Backend} (P4, P5, P6, P7, P8) & D & \textit{API Flooding} ke P4 (Proxy) atau P8 (IOTA Client). & Layanan \textit{backend} lumpuh. \\
    \hline
    T16 & Proses \textit{Backend} (P4, P5, P6, P7, P8) & E & Memanipulasi logika CapBAC saat verifikasi. & Eskalasi hak akses permanen. \\
    \hline
    T17 & Penyimpanan Internal (D1, D3) & T & Modifikasi \textit{nonce} di Redis / manipulasi \textit{Key Store}. & \textit{Replay attack} berhasil. \\
    \hline
    T18 & Penyimpanan Internal (D1, D3) & I & Basis data memori diakses pihak ketiga. & Materi kunci/token dicuri. \\
    \hline
    T19 & Penyimpanan Internal (D1, D3) & D & Menghapus paksa isi \textit{cache} atau direktori \textit{key}. & Gagal dekripsi/otentikasi. \\
    \hline
    T20 & Penyimpanan Publik (D2, D4) & T & Modifikasi \textit{ciphertext} di IPFS / Mayoritas \textit{node attack}. & Integritas rekam medis hancur. \\
    \hline
    T21 & Penyimpanan Publik (D2, D4) & R & (Khusus D4) Kehilangan log akses \textit{immutable}. & Jejak riwayat tak terlacak. \\
    \hline
    T22 & Penyimpanan Publik (D2, D4) & I & Analisis relasi lalu lintas metadata di publik. & Privasi institusi/pasien bocor. \\
    \hline
    T23 & Penyimpanan Publik (D2, D4) & D & \textit{Spam} \textit{junk data} / transaksi palsu ke \textit{node}. & Penarikan RME melambat. \\
    \hline
    T24 & Seluruh Aliran Data (F1 - F29) & T & \textit{Man-in-the-Middle} (MitM) mengubah paket jaringan. & Perintah transit termanipulasi. \\
    \hline
    T25 & Seluruh Aliran Data (F1 - F29) & I & \textit{Sniffing} pada jalur tanpa TLS/enkripsi memadai. & \textit{Payload} jaringan disadap. \\
    \hline
    T26 & Seluruh Aliran Data (F1 - F29) & D & Pemutusan koneksi / \textit{flooding} jalur komunikasi. & Sinkronisasi data terhenti. \\
    \hline
\end{xltabular}